\section{Inleiding}
Het programma is een toepassing van \textit{Decision tree learning}. \cite{wiki:dtl}
\section{Teamwerk}
Één programma schrijven met meerdere programmeurs is niet simpel. Daarom gebruiken wij twee belangrijke tools. Met deze tools konden wij ons meer focussen op het programma zelf en minder op het oplossen van randproblemen.

\subsection{Git}
De eerste tool die wij gebruiken is Git.
Git is een \textit{Version Control System} (VCS) oftewel een versiebeheersysteem. De oorspronkelijke ontwikkelaar van git is Linus Torvalds in 2005.\cite{git:init} Ondertussen werken er al veel meer ontwikkelaars aan.

In een zogenoemde \textit{Git repository} worden alle verandering van code onthouden. Deze veranderingen slaagt men als programmeur op per nieuwe revisie. Deze revisies worden allemaal bijgehouden op chronologische volgorde. Zo kan men teruggaan naar een vorige revisie en deze vergelijken met de nieuwere revisie. Zo kan men de oorzaak van een fout sneller opsporen.

\subsection{Github}
De website \textit{www.github.com}, oftewel Github, is de tweede tool.
Deze website is een online platform voor versiebeheer via Git. Via Github kunnen programmeurs hun \textit{repository} online opslaan.\cite{git:hello_world} Zo kunnen zij altijd en overal aan hun code en deze ook snel delen met andere programmeurs.

\begin{figure}[ht]
    \centering
    \includegraphics[width=0.4\textwidth]{illustraties/Octocat}
    \caption{Octocat, de mascotte van GitHub}
    \cite{git:github}
    \label{fig:octocat}
\end{figure}

\section{Boomstructuur}
\subsection{Tree-object}
\subsection{TreeNode-object}

\section{JSON}
De regressieboom wordt gegenereerd in python. Deze wordt dan gecodeerd en weggeschreven in JSON-formaat. Voor het inlezen van data in JSON-formaat bestaan al goede bibliotheken. Toch hebben we gekozen om dit zelf te programmeren. Wij waren namelijk ook intrinsiek gemotiveerd uit te zoeken hoe zo'n parser zou kunnen werken.

\begin{figure}[ht]
    \centering
    \lstinputlisting[language=JavaScript]{code/rules_d2.json}
    \caption{De boomstructuur in JSON-formaat}
    \label{fig:json_tree}
\end{figure}

Zoals te zien op Figuur \ref{fig:json_tree}, heeft elke node twee zogenoemde \textit{key-value pairs} of sleutel-waardeparen: \lstinline[language=JavaScript]{"key": "value"}. 
Deze zijn telkens gescheiden door komma's. Deze paren staan in een onbepaalde volgorde. Tijdens het lezen moeten we dus controleren waar elke sleutel voor staat.

\subsection{Sleutel: name}
In deze sleutel kunnen twee dingen staan.
\subsection{Sleutel: children}

\section{Inladen van de structuur}

\section{Regex}
De beslissing of de prijs wordt opgeslagen in een string. Om deze data snel en foutloos uit deze string te halen gebruiken we regular expressions oftewel regex. Een zogenoemde regular expression is een reeks karakters die samen een patroon definiëren. Zo'n patroon wordt gebruikt om bepaalde dingen te zoeken, om speciefieke dingen te veranderen of om inputs te controleren.\cite{wiki:regex}

\subsection{Prijs}
\begin{lstlisting}[language=JavaScript]
{
    "name": "267 of Price",
    "children": null
}
\end{lstlisting}
\subsection{Beslissing}

\begin{lstlisting}[language=JavaScript]
{
    "name": "Model_B3 > 0.5",
    "children": [{...}, {...}]
}
\end{lstlisting}

\section{Testen}
In dit deel wordt getest hoe goed de pijsbepaling gebeurt bij bepaalde boomdieptes. Er wordt ook getest wat de invloed is van de boomdiepte op de tijd die het programma nodig heeft om de boom in te laden en om de prijs te bepalen. 

Om een zo accuraat mogelijk resultaat te bekomen, hebben we elke boom tienduizend, \(10^4\), keer ingeladen. Ook voor elke boom hebben we tien miljoen, \(10^7\), orgels geprijsd. De test duurde in totaal \(8.15613\) seconden. De resultaten zijn te zien in Tabel \ref{fig:test_result}.

\begin{table}[ht]
    \small
    \centering
    \begin{tabular}{|c|r|r|r|r|l|}
        \hline
        \bfseries Depth & \bfseries Load time & \bfseries Estimate time & \bfseries Total error & \bfseries Abs. error & \bfseries Rel. error\\
        \hline
        \csvreader[head to column names, late after line=\\]{data.csv}{}%
        {\depth & \load s & \est s & \total & \abs & \rel \%}
        \hline
    \end{tabular}
    \caption{Testresultaten}
    \label{tab:test_results}
\end{table}

De inlaadtijd van de boom is duidelijk afhankelijk van de boomdiepte. Bij een diepere boom zien we een duidelijke stijging in inlaadtijd. Met de grote-O-notatie is dit een complexiteit van $ O({2^n}) $. In dit geval is het geen worst case scenario omdat de boom niet overal vertakt tot beneden.

Bij het bepalen van de prijs zien we ook dat het langer duurt bij een diepere boom.

Bij het testen gebruiken we alleen orgels waarvan we de prijs kennen. Het programma maakt dan de schatting van de prijs. Deze prijs wordt dan vergeleken met echte prijs. Bij het nemen van de absolutewaarde van de fout en hier het gemmiddelde van te nemen over alle orgels krijgen we de absolute fout die te vinden is in de tabel. Het is heel dijdelijk dat deze fout kleiner wordt als de boom dieper wordt. Dit komt natuurlijk door dat er meer beslissing worden genomen en dus een betere prijsschating kan worden gedaan. Als we dan naar de procentuele fout gaan kijken, is deze bij een diepte van 5 heel klein. Daaruit kunnen we besluiten dat we bij een diepte van 5 een precise bepaling hebben van de pijs.

\begin{figure}[ht]
    \centering
    \begin{tikzpicture}
        \begin{axis}[
            width=0.7\textwidth,
            xlabel={Boom diepte },
            ylabel={Gem. ladingstijd (s)}]
            \addplot table [col sep=comma, x=depth , y=load] {data.csv};
        \end{axis}
    \end{tikzpicture}

    \begin{tikzpicture}
        \begin{axis}[
            width=0.7\textwidth,
            xlabel={Boom diepte },
            ylabel={Gem. uitvoertijd (s)}]
            \addplot table [ col sep=comma, x=depth, y=est] {data.csv};
        \end{axis}
    \end{tikzpicture}

    \caption{Grafieken Testresultaten}
    \label{fig:test_result}
\end{figure}

\subsection{Testomgeving}
In Tabel \ref{tab:pc-specs} staan de specificaties van de gebruikte computer. Van alle computers waarop we de test konden uitvoeren, was deze computer de snelste.

\begin{table}[ht]
    \centering
    \begin{tabular}{|l|l|}
        \hline
        \multicolumn{2}{|l|}{Software} \\
        \hline
        VS Build configuration  & Release x64 \\
        Besturingssysteem       & Windows 10 Pro 64-bit (10.0, Build 17763) \\
        \hline
        \hline
        \multicolumn{2}{|l|}{Hardware} \\
        \hline
        Processor               & Intel Core i7-7700K CPU @ 4.60 GHz OC (8 CPU's) \\
        Geheugen (RAM)          & 2x8.0 GB @ 2933 MHz \\
        Grafische kaart         & NVIDIA GeForce GTX 1080 Ti \\
        \hline
    \end{tabular}
    \caption{Specificaties van de gebruikte computer}
    \label{tab:pc-specs}
\end{table}


% We fixen de positie van deze figuur nog wel
